\section{Higher-Order Collision Geometry}
\label{app:higher-order-collision-geometry}

The main development uses the pairwise collision kernel $k_\theta\equiv k_{\theta,2}$, but the same event has an integer-order extension.
For integer $\alpha\ge2$ and conditionally independent draws $Z_m\sim q_\theta(\cdot\mid x_m)$,
\begin{equation}
k_{\theta,\alpha}(x_1,\ldots,x_\alpha)
=\sum_z\prod_{m=1}^{\alpha}q_\theta(z\mid x_m)
=\Pr[Z_1=\cdots=Z_\alpha\mid x_1,\ldots,x_\alpha].
\label{eq:higher-order-collision}
\end{equation}
A sequence converging to deterministic assignments gives the indicator that every element of the tuple belongs to the same hard class.
At order two, $k_{\theta,2}(x,x')=k_\theta(x,x')$, recovering \cref{eq:collision-kernel}.
Collision arity is a property of this representation statistic; it need not equal the arity of a source relation.

For even order $\alpha=2m$, define the elementwise product feature
\begin{equation}
\phi_m(x_1,\ldots,x_m)
=q_\theta(\cdot\mid x_1)\odot\cdots\odot q_\theta(\cdot\mid x_m).
\end{equation}
Then
\begin{equation}
k_{\theta,2m}(x_1,\ldots,x_{2m})
=\langle\phi_m(x_1,\ldots,x_m),\phi_m(x_{m+1},\ldots,x_{2m})\rangle.
\label{eq:even-order-kernel}
\end{equation}
It is an ordinary positive-semidefinite kernel on $m$-tuples.
Odd order $2m+1$ gives the analogous inner product of $\phi_m$ and $\phi_{m+1}$, naturally a cross-affinity between tuple spaces of different arities rather than a symmetric kernel on individual inputs.

\paragraph{Triplet-structured relational requirements}
A learning objective may involve three inputs while being composed entirely from pairwise collision geometry.
Writing $q_a:=q_\theta(\cdot\mid x_a)$, $q_p:=q_\theta(\cdot\mid x_p)$, and $q_n:=q_\theta(\cdot\mid x_n)$, and reusing the spherical map $u(q)$ from \cref{eq:sphere-map}, one triplet-structured angular objective is
\begin{equation}
\mathcal L_{\mathrm{trip}}
=
\left[
\zeta-u(q_a)^\top u(q_p)+u(q_a)^\top u(q_n)
\right]_+,
\qquad \zeta>0.
\end{equation}
This objective compares pairwise normalized collisions from \cref{eq:normalized-collision}; it is not a third-order collision statistic.
Angular comparison of high-dimensional graph representations also appears in graph summarization through cosine-distance spherical clustering~\citep{tsalouchidou2020scalable}.
There it acts on graph-derived connectivity-history vectors, whereas here it acts on categorical assignment distributions, for which $1-u(q)^\top u(q')=1-\cos\vartheta$ is cosine distance and the underlying unnormalized inner product has a collision-probability interpretation.

A genuine third-order collision instead requires all three independently encoded inputs to receive the same codeword:
\begin{equation}
k_{\theta,3}(x_a,x_p,x_n)
=
\sum_z q_\theta(z\mid x_a)q_\theta(z\mid x_p)q_\theta(z\mid x_n)
=
\Pr[Z_a=Z_p=Z_n].
\end{equation}
The related anchor--favored--disfavored event is
\begin{equation}
\Pr[Z_a=Z_p\ne Z_n]
=
\sum_z q_\theta(z\mid x_a)q_\theta(z\mid x_p)
\bigl(1-q_\theta(z\mid x_n)\bigr)
=
k_\theta(x_a,x_p)-k_{\theta,3}(x_a,x_p,x_n),
\end{equation}
with hard-code limit
\[
\ind[c_\theta(x_a)=c_\theta(x_p),\;c_\theta(x_a)\ne c_\theta(x_n)].
\]
Thus source-relation arity, learning-objective arity, and collision-event order are distinct modelling choices.

For iid inputs $X_1,\ldots,X_\alpha\sim\mu$,
\begin{equation}
\E[k_{\theta,\alpha}(X_1,\ldots,X_\alpha)]
=\sum_zq_Z(z)^\alpha.
\label{eq:higher-order-population-collision}
\end{equation}
The factorization follows by independence of the $\alpha$ inputs.
At $\alpha=2$, this identity recovers \cref{eq:population-collision}.
Consequently, integer-order R\'enyi entropy~\citep{renyi1961measures} satisfies
\begin{equation}
H_\alpha(q_Z)
=\frac1{1-\alpha}\log\sum_zq_Z(z)^\alpha
=-\frac1{\alpha-1}\log\E[k_{\theta,\alpha}].
\label{eq:renyi-alpha-collision}
\end{equation}
The case $\alpha=2$ recovers \cref{eq:collision-entropy}.
The literal $\alpha$-way same-code event is defined here for integer $\alpha\ge2$, whereas the R\'enyi family extends to appropriate real orders under its usual conventions.
Shannon entropy arises through the continuous limit $\alpha\to1$, not from a nontrivial one-sample collision event.

For a fixed positive reference $r$, define
\begin{equation}
a_{\alpha,r}(x_1,\ldots,x_\alpha)
=\sum_zr(z)^{1-\alpha}\prod_{m=1}^{\alpha}q_\theta(z\mid x_m).
\end{equation}
$a_{2,r}(x,x')=a_r(x,x')$ recovers the reference-weighted affinity in \cref{eq:reference-affinity}.
Its expectation is $\sum_zq_Z(z)^\alpha r(z)^{1-\alpha}$, giving
\begin{equation}
D_\alpha(q_Z\|r)
=\frac1{\alpha-1}\log\E[a_{\alpha,r}].
\label{eq:renyi-alpha-reference-affinity}
\end{equation}
For a uniform reference,
\begin{equation}
D_\alpha(q_Z\|U_K)=\log K-H_\alpha(q_Z).
\label{eq:uniform-renyi-alpha}
\end{equation}
At $\alpha=2$, this reduces to \cref{eq:uniform-renyi}.
This includes the usual limiting and support conventions at orders one and zero.
At fixed $K$, minimizing this divergence as an organization penalty encourages larger occupancy entropy, whereas adding $+\lambda H_\alpha(q_Z)$ to a minimized objective encourages concentration.

\subsection{R\'enyi entropy in spherical coordinates}
\label{app:renyi-spherical-coordinates}

The integer-event interpretation and the real-order entropy family should be distinguished.
For $\alpha>0$, $\alpha\ne1$, let $u=u(q)$ be the spherical map from \cref{eq:sphere-map}.
Substitution of $q_i=u_i/\|u\|_1$ yields
\begin{equation}
H_\alpha(q)
=\frac1{1-\alpha}\left[\log\sum_i u_i^\alpha-\alpha\log\|u\|_1\right].
\label{eq:renyi-spherical}
\end{equation}
At order two, $\sum_i u_i^2=1$ gives $H_2(q)=2\log\|u\|_1$.
Taking the order-one limit gives \cref{eq:shannon-spherical}.
All these quantities can be expressed in the same coordinates, but higher power sums generally depend on more than one ordinary spherical angle.

\section{Conditional Collision and Dependence}
\label{app:conditional-collision-dependence}

The aggregate population collision and $H_2(q_Z)$ developed in \cref{eq:population-collision,eq:collision-entropy} describe how a finite code is occupied and organized across a population, but they do not by themselves determine how strongly the sampled code depends on its input.
In particular, broad aggregate occupancy can arise even from an input-independent encoder.
This appendix records two collision-related dependence functionals that make this distinction explicit: one compares conditional self-collision with aggregate population collision, whereas the other uses R\'enyi divergence of the induced input--code joint distribution from its independent product.
It also distinguishes input--code dependence from dependence between relational endpoints, for which collision alignment observes only the same-codeword event.

Let $X,X'\overset{\mathrm{iid}}{\sim}\mu$.
Using the collision kernel and self-collision in \cref{eq:collision-kernel,eq:self-collision}, together with the aggregate code distribution and population collision identity in \cref{eq:aggregate-code,eq:population-collision}, define
\begin{equation}
\kappa_{\mathrm{self}}
:=\E_X[k_\theta(X,X)]
=\E_X\sum_zq_\theta(z\mid X)^2,
\qquad
\kappa_{\mathrm{pop}}
:=\E_{X,X'\overset{\mathrm{iid}}{\sim}\mu}[k_\theta(X,X')]
=\sum_zq_Z(z)^2.
\end{equation}
The conditional collision-entropy functional is
\begin{equation}
H_2^{\mathrm{coll}}(Z\mid X)=-\log \kappa_{\mathrm{self}},
\end{equation}
Because the logarithm is taken after averaging conditional self-collision, this functional is generally not equal to $\E_X[H_2(q_\theta(\cdot\mid X))]$.
Using the collision entropy $H_2(q_Z)$ from \cref{eq:collision-entropy}, a corresponding difference is
\begin{equation}
I_2^{\mathrm{coll}}(X;Z)
=H_2(q_Z)-H_2^{\mathrm{coll}}(Z\mid X)
=\log\frac{\kappa_{\mathrm{self}}}{\kappa_{\mathrm{pop}}}.
\label{eq:collision-information}
\end{equation}
It is nonnegative because
\begin{equation}
\kappa_{\mathrm{self}}-\kappa_{\mathrm{pop}}
=\sum_z\operatorname{Var}_\mu(q_\theta(z\mid X))\ge0.
\end{equation}
It vanishes exactly when the assignment vector is constant almost surely, equivalently when the sampled code is independent of the input under this population model.
For deterministic assignments, $\kappa_{\mathrm{self}}=1$ and the difference equals $H_2(q_Z)$.
No general data-processing or unique R\'enyi-mutual-information characterization is asserted for this entropy difference.

Let $P_X$ denote the distribution of $X$, so $P_X=\mu$, and let $P_{XZ}$ be the joint distribution induced by $P_X$ and $q_\theta(\cdot\mid X)$; its $Z$-marginal is the aggregate distribution $q_Z$ from \cref{eq:aggregate-code}.
Using the inverse-mass affinity of \cref{eq:qz-affinity}, a divergence-based quantity has a different form:
\begin{equation}
I_2^D(X;Z)
=D_2\!\left(P_{XZ}\,\middle\|\,P_X\otimes q_Z\right)
=\log\E_X[a_{q_Z}(X,X)].
\label{eq:divergence-renyi-information-affinity}
\end{equation}
The expansion follows on the active support from the density ratio $q_\theta(z\mid x)/q_Z(z)$.
For deterministic assignments, its exponential is the number of active words, so $I_2^D=R_0=\log|\operatorname{supp}(q_Z)|$ from \cref{eq:finite-state-rate-hierarchy} rather than $H_2(q_Z)$.
The two functionals illustrate why an unqualified reference to R\'enyi mutual information would be ambiguous here.

There is a further distinction between input--code dependence and dependence of relational endpoints.
For $(\mathsf V,\mathsf V')\sim P_{\mathrm{pair}}$ as in \cref{eq:generic-relational-objective}, with endpoints encoded independently from $q_\theta(\cdot\mid x_{\mathsf V})$ and $q_\theta(\cdot\mid x_{\mathsf V'})$, the full joint distribution of $(Z,Z')$ can contain information beyond the equality event.
Collision alignment measures only
\begin{equation}
\Pr[Z=Z']=\E_{(\mathsf V,\mathsf V')\sim P_{\mathrm{pair}}}[k_\theta(x_{\mathsf V},x_{\mathsf V'})].
\end{equation}
It is an equality-specific statistic, not the full mutual information between the endpoints.

\section{Additional Monotone Content Examples}
\label{app:path-walk-content}

The main text defines content-based relational fidelity by choosing a nonnegative structural-content functional $\mathcal I$ that is monotone under an admissible simplification and measuring the resulting content loss through \cref{eq:relational-content-monotone,eq:relational-content-distortion}.
The direct-edge realization in \cref{eq:edge-distortion} uses retained edge weight, but the same construction can emphasize structure induced by several relations rather than individual edge entries.
This appendix records shortest-path, walk, and natural-connectivity examples to make that latitude concrete; they are illustrative alternatives and are not evaluated in the experiments.

For a graph $G=(V,E)$ with positive edge lengths and an edge-deletion simplification $\widetilde G$ on the same carrier, let $d_G(i,j)$ be shortest-path distance, using infinity when no path exists.
If $\varphi:[0,\infty]\to\R_{\ge0}$ is nonincreasing and $\varphi(\infty)=0$, then
\begin{equation}
\mathcal I_\varphi(G)=\sum_{i<j}\varphi(d_G(i,j))
\end{equation}
is monotone under edge deletion.
Deleting edges cannot shorten a path.
For example, with $\varphi(d)=e^{-d/\tau}$ and $\tau>0$, the corresponding specialization of the generic content distortion in \cref{eq:relational-content-distortion} is
\begin{equation}
D_\varphi(G,\widetilde G)
=\sum_{i<j}\left[e^{-d_G(i,j)/\tau}-e^{-d_{\widetilde G}(i,j)/\tau}\right].
\end{equation}
This values the geometry supported by paths rather than only the removed edge entries.
For weight attenuation, a relationship between weights and lengths would additionally need to be specified.

For the nonnegative weighted adjacency $W$ used in \cref{sec:edge-cut-distortion}, and an admissible attenuation with reconstructed adjacency $\widetilde W$, finite walk content can be written
\begin{equation}
\mathcal I_{\mathrm{walk}}(G)
=\sum_{\ell=1}^{L_0}\alpha_\ell\mathbf1^\top W^\ell\mathbf1,
\qquad\alpha_\ell\ge0.
\end{equation}
Each term counts weighted walk mass of one length.
If $0\le\widetilde W\le W$ entrywise, induction using nonnegative matrix multiplication gives $0\le\widetilde W^\ell\le W^\ell$ entrywise.
Hence the content is monotone under attenuation.
The truncation $L_0$ is used to avoid confusing the walk length limit with the graph Laplacian $L$.

Natural connectivity is an established spectral example~\citep{wu2010natural}.
For a finite undirected unweighted graph $G=(V,E)$ with adjacency $A$ and $n=|V|$,
\begin{equation}
\mathcal I_{\mathrm{NC}}(G)
=\log\left(\frac{\operatorname{tr}(e^A)}{n}\right)
=\log\left(\frac1n\sum_r e^{\lambda_r(A)}\right).
\end{equation}
The expansion $\operatorname{tr}(e^A)=\sum_{\ell\ge0}\operatorname{tr}(A^\ell)/\ell!$ connects it to closed walks.
It is zero for the empty graph and strictly increases when an edge is added in this standard setting.
Natural connectivity supplies an established structural-content functional for this construction.

Direct edge content in \cref{eq:edge-distortion} treats each primitive relation according to its weight, shortest-path content responds to induced connectivity geometry, finite walk content measures multistep structure, and natural connectivity provides an all-orders closed-walk spectral example.
A fidelity criterion may combine monotone contents with nonnegative coefficients, with its interpretation tied to the declared combination of relational properties.

\section{Reported Metrics and Numerical Conventions}
\label{app:metric-definitions}

This appendix collects the definitions of recurring metrics reported in the experiments and the aggregation conventions needed to interpret them.
Metrics defined earlier retain the same notation and are referenced back to their original definitions.

\subsection{Hard finite codes}

For $N$ evaluated hard codewords $c(i)\in\mathcal Z$, let
\begin{equation}
\widehat q_Z(z)
=\frac1N\sum_{i=1}^N\ind[c(i)=z].
\end{equation}
The distribution $\widehat q_Z$ is the empirical codeword distribution over the evaluated hard tokens and is the empirical counterpart of the aggregate distribution $q_Z$ introduced in \cref{eq:aggregate-code}.

Aggregate hard collision is the probability that two independent draws from $\widehat q_Z$ receive the same codeword:
\begin{equation}
\widehat\kappa_{\mathrm{pop}}
=\sum_z\widehat q_Z(z)^2,\qquad
\widehat H_2
=-\log\widehat\kappa_{\mathrm{pop}},\qquad
K_{\mathrm{eff}}
=\widehat\kappa_{\mathrm{pop}}^{-1}.
\end{equation}
Here $\widehat\kappa_{\mathrm{pop}}$ is the empirical same-code collision, $\widehat H_2$ is the hard-code collision entropy, and $K_{\mathrm{eff}}$ is the collision-effective number of codewords.
They are the empirical counterparts of the population quantities in \cref{eq:population-collision,eq:collision-entropy,eq:effective-code-count}.
The reported number of active hard words is $|\{z:\widehat q_Z(z)>0\}|$.

For the $32\times32$ latent grid on $256\times256$ images,
\[
\mathrm{bpp}
=\frac{32^2b}{256^2}
=\frac{b}{64}.
\]
This is the nominal raw latent bit capacity used in Study~4 before entropy coding.

\subsection{Graph partitions}

For a hard assignment $c$, use the classes $V_z=\{i:c(i)=z\}$ from \cref{prop:normalized-cut}, with
\[
\operatorname{vol}(G)=\sum_i d_i,\qquad
\operatorname{vol}(V_z)=\sum_{i\in V_z}d_i.
\]
Using directed symmetric edge incidences,
\[
\operatorname{assoc}(V_z,V_z)
=\sum_{(i,j)\in\vec E:i,j\in V_z}W_{ij}.
\]
The reported hard fixed-$K$ normalized cut is the specialization in \cref{eq:hard-ncut}:
\[
\operatorname{NAssoc}(G)
=\sum_{z:\operatorname{vol}(V_z)>0}
\frac{\operatorname{assoc}(V_z,V_z)}{\operatorname{vol}(V_z)},\qquad
\operatorname{Ncut}(G)
=K-\operatorname{NAssoc}(G).
\]
Under the fixed-$K$ convention, an empty class contributes zero association and therefore one unit to $\operatorname{Ncut}(G)$.
The reported largest hard volume fraction is $\max_z\operatorname{vol}(V_z)/\operatorname{vol}(G)$, and active hard classes are those with positive volume.
The graph-fidelity values reported in Study~2 are the hard specializations of $D_E$, $D_F$, $D_C$, and $D_{H_2}$ defined in \cref{eq:normalized-edge-distortion,eq:source-normalized-dirichlet-distortion,eq:source-collision-edge-distortion,eq:entropy-probe-distortion}.
The post-hoc $D_F$ reported in Study~3 uses the hard-partition form in \cref{eq:hard-partition-dirichlet-distortion}.

For graph-local assignments, $\bar q_G(z)$ is the degree-weighted occupancy defined in \cref{eq:graph-aggregate-code}.
The reported organization metrics are
\[
H_2(\bar q_G)
=-\log\sum_z\bar q_G(z)^2,\qquad
K_{\mathrm{eff}}
=\left(\sum_z\bar q_G(z)^2\right)^{-1},
\]
\[
D_2(\bar q_G\|U_K)
=\log K-H_2(\bar q_G)
=\log\left(K\sum_z\bar q_G(z)^2\right).
\]
Following \cref{eq:collision-entropy,eq:effective-code-count,eq:uniform-renyi}, $H_2(\bar q_G)$ is graph-local collision occupancy, $K_{\mathrm{eff}}$ is the corresponding collision-effective number of classes, and $D_2(\bar q_G\|U_K)$ measures deviation from uniform graph-local occupancy.
These quantities are computed graph-locally before averaging across evaluated graphs.
Consequently, the mean of $K_{\mathrm{eff}}$ is not generally the exponential of the mean $H_2$, and mean $D_2$ cannot be converted to a mean effective count by exponentiation.
The learned MalNet results report the mean and sample standard deviation across five seeds, while the spectral reference is deterministic for the stated solver configuration.

\subsection{Teacher relations and reconstruction}

As introduced in \cref{sec:flowers102-teacher-experiment}, the favored and disfavored weighted hard-code collision rates are
\begin{equation}
\begin{aligned}
\widehat\kappa_+&=\frac{\sum_{ij}S_{+,ij}\ind[c(i)=c(j)]}{\sum_{ij}S_{+,ij}},\\
\widehat\kappa_-&=\frac{\sum_{ij}S_{-,ij}\ind[c(i)=c(j)]}{\sum_{ij}S_{-,ij}}.
\end{aligned}
\end{equation}
The pooled empirical hard-code collision is $\widehat\kappa_{\mathrm{pop}}$ as defined above.
Favored-pair enrichment is $\widehat\kappa_+/\widehat\kappa_{\mathrm{pop}}$, while disfavored-pair suppression is $1-\widehat\kappa_-/\widehat\kappa_{\mathrm{pop}}$.
The reported favored and disfavored teacher-cosine summaries are the corresponding relation-weighted means:
\[
\frac{\sum_{ij}S_{\pm,ij}\,e_i^\top e_j}
{\sum_{ij}S_{\pm,ij}}.
\]
Validation and test teacher loss use the alignment and relational-separation terms of \cref{eq:dino-teacher-loss}, with their numerators aggregated using the corresponding $S_+$ and $S_-$ weight masses before forming the final average.

For a unit-range RGB image $x$ and reconstruction $\widehat x$,
\[
\operatorname{MSE}(x,\widehat x)
=\frac{1}{3HW}\|x-\widehat x\|_2^2.
\]
Reported MSE averages this squared pixel error over the evaluated images.
For unit-range images,
\[
\operatorname{PSNR}
=-10\log_{10}(\operatorname{MSE}),
\]
so higher PSNR corresponds to lower MSE.
MS-SSIM is the standard multiscale structural similarity diagnostic of \citet{wang2003multiscale}.

\section{Synthetic Study Details}
\label{app:synthetic-details}

The synthetic study uses one fixed finite dataset and evaluates every point pair exactly.
Two copies of the same encoder are initialized identically and optimized using the centroid and inverse-mass-weighted pairwise forms, respectively.
The separate decoder calculation holds the encoder assignments fixed.
These checks verify the finite-sample centroid--pairwise identity in \cref{eq:finite-soft-pairwise} and the decoder decomposition in \cref{eq:decoder-centroid-decomposition}.
If $g_{\mathrm{cent}}$ and $g_{\mathrm{pair}}$ are the flattened encoder gradients of the centroid and pairwise objectives, their reported gradient agreement uses cosine similarity and relative error
\[
\frac{\|g_{\mathrm{cent}}-g_{\mathrm{pair}}\|_2}
{\|g_{\mathrm{cent}}\|_2}.
\]

\begin{table}[htbp]
\centering\small
\begin{tabularx}{\linewidth}{Xr}
\toprule
Diagnostic & Observed value \\
\midrule
Maximum $|D_{\mathrm{cent}}-D_{\mathrm{pair}}|$, centroid-trained trajectory & $9.54\times10^{-7}$ \\
Maximum $|D_{\mathrm{cent}}-D_{\mathrm{pair}}|$, pair-trained trajectory & $9.54\times10^{-7}$ \\
Larger trajectory mean of $|D_{\mathrm{cent}}-D_{\mathrm{pair}}|$ & $4.20\times10^{-8}$ \\
Random encoder states: maximum $|D_{\mathrm{cent}}-D_{\mathrm{pair}}|$ & $1.91\times10^{-6}$ \\
Random encoder states: mean $|D_{\mathrm{cent}}-D_{\mathrm{pair}}|$ & $2.64\times10^{-7}$ \\
Initial gradient cosine / relative error & $1.000\,/\,1.54\times10^{-7}$ \\
Final gradient cosine / largest relative error across final models & $1.000\,/\,1.35\times10^{-6}$ \\
Final parameter RMS difference between the centroid- and pairwise-trained encoders & $3.56\times10^{-8}$ \\
Random states: mean $|D_{\mathrm{cent}}-KD_{\mathrm{raw}}|$ & $4.79$ \\
Random states: maximum $|D_{\mathrm{cent}}-KD_{\mathrm{raw}}|$ & $21.69$ \\
Decoder excess term in \cref{eq:decoder-centroid-decomposition}, steps $0\to200$ & $21.05\to8.15\times10^{-5}$ \\
Maximum decoder decomposition residual & $1.91\times10^{-6}$ \\
\bottomrule
\end{tabularx}
\caption{Synthetic numerical diagnostics. Values near $10^{-6}$ reflect single-precision arithmetic.}
\label{tab:synthetic-collision-centroid}
\end{table}

\begin{table}[htbp]
\centering\small
\begin{tabularx}{\linewidth}{>{\raggedright\arraybackslash}p{0.28\linewidth}X}
\toprule
Component & Setting \\
\midrule
Data & 768 points in $\R^2$ from a six-component Gaussian mixture \\
Mixture weights & $(0.30,0.22,0.17,0.13,0.10,0.08)$ \\
Realized component counts & $(230,169,131,100,77,61)$ \\
Encoder & MLP with two hidden layers, width 32, $\tanh$ nonlinearities \\
Code distribution & Exact enumeration of $2^3=8$ factorized Bernoulli words \\
Code temperature & 0.7 \\
Encoder optimization & SGD, learning rate 0.02, momentum 0.9, 300 steps \\
Random-state check & 128 random affine-logit maps; weight/bias scale log-uniform in $[0.25,4.0]$ \\
Decoder check & Learned $8\times2$ reproduction codebook, SGD learning rate 0.15, 200 steps \\
Decoder initialization & Gaussian reproduction vectors with scale 3.0 \\
Seed & 1337 \\
\bottomrule
\end{tabularx}
\caption{Settings for Study 1.}
\label{tab:synthetic-settings}
\end{table}

\section{Transductive Relational-Fidelity Study Details}
\label{app:transductive-relational-distortion-details}

The transductive study optimizes a separate logit matrix for each source graph, fidelity, and organization weight.
The two MalNet collections use the same 20 cleaned MalNet-Tiny test topologies.
The weighted variant assigns deterministic positive lognormal weights, with scale $\sigma=0.75$, mean-normalized within each graph, and global seed 91337.
COLLAB and PROTEINS each contribute 20 cleaned largest connected components from TUDataset.
All labels are removed or ignored.

The solver uses ten random initializations and one collapse-biased initialization, retaining the best target soft objective encountered in each optimization.
After these ordinary restarts, each target is warm-started from the other criteria's selected logits; if this improves any target, one closure pass is run from the updated solutions.
Final selection uses the target soft objective across all candidates.

\begin{table}[htbp]
\centering\small
\begin{tabularx}{\linewidth}{>{\raggedright\arraybackslash}p{0.28\linewidth}X}
\toprule
Component & Setting \\
\midrule
Collections & 20 MalNet-Tiny test LCCs; the same 20 weighted topologies; 20 COLLAB; 20 PROTEINS \\
Base graph-selection seed & 1337; deterministic collection-specific offsets for COLLAB and PROTEINS \\
Minimum source size & 512 nodes for MalNet; 16 nodes for TUDataset, after LCC extraction \\
Training fidelities & $D_E$, $D_F$, $D_C$, and $D_{H_2}$ where defined \\
Entropy-field exclusions & Omit $D_{H_2}$ when its source Dirichlet denominator is zero; two selected COLLAB graphs are excluded from this criterion \\
Representation & Free graph-local $n\times8$ categorical logits; no learned encoder \\
Assignment temperature & $\tau=1.0$ \\
Organization weights & $0$, $0.03$, $0.05$, $0.07$, $0.08$, $0.09$, $0.10$, $0.11$, $0.12$, $0.14$, $0.16$, $0.18$, $0.20$, $0.30$, $0.50$ \\
Objective & $D_\rho(q)+\lambda_{\mathrm{org}}D_2(\bar q_G\|U_8)$ \\
Optimizer & Adam, learning rate 0.05, 600 steps, no scheduler \\
Random initialization & Gaussian logits with standard deviation $10^{-2}$ \\
Random restart seeds & 1337, 2024, 31415, 2718, 1618, 9001, 42, 73, 101, 211 \\
Collapse-biased restart & Initial logit bias 4.0 toward one state \\
Warm starts & Cross-criterion pass from the other selected solutions; one closure pass if the first pass improves any target \\
Selection & Best target soft objective across ordinary restarts and cross-criterion candidates \\
Diagnostics & Hard/soft cross-evaluation, source-importance similarities, and 256 size-balanced random partitions per graph \\
Random-partition seed & $1337+\text{graph index}$ \\
\bottomrule
\end{tabularx}
\caption{Settings for Study 2.}
\label{tab:transductive-relational-distortion-settings}
\end{table}

\subsection{Matched occupancy and cross-distortion transfer}

Hard partitions are compared at target $H_2(\bar q_G)$ values from 0.25 to 2.0 nats.
Each selected solution must be within 0.10 nats of the target, and the compared pair must be within 0.10 nats of one another.
Each criterion selects its solution independently from the available sweep before this matching rule is applied.
At least five matched graphs are required to display a point in the agreement figures.
Consequently, different source/criterion comparisons can use different subsets, and missing targets reflect available coverage rather than an interpolation through unobserved solutions.

\Cref{tab:transductive-rate-matched-agreement} preserves the detailed comparison at 2.0 nats.
Let $c_A$ and $c_B$ denote the partitions trained under fidelities $A$ and $B$, respectively.
The directional transfer differences are
\[
\Delta_{A\to B}=D_B(c_A)-D_B(c_B),
\qquad
\Delta_{B\to A}=D_A(c_B)-D_A(c_A).
\]
Each difference is measured on the destination fidelity's scale.
Because matching is approximate and selection uses the soft training objective, these descriptive differences can be slightly negative.

\begin{table}[htbp]
\centering\small
\setlength{\tabcolsep}{4pt}
\begin{tabular}{llrrrr}
\toprule
Source & Pair & $n$ & Mean ARI & $\Delta_{A\to B}$ & $\Delta_{B\to A}$ \\
\midrule
MalNet unweighted & $D_F$ / $D_C$ & 20 & 0.781 & -0.0025 & 0.0060 \\
MalNet unweighted & $D_F$ / $D_{H_2}$ & 20 & 0.166 & 0.0919 & 0.0698 \\
MalNet unweighted & $D_C$ / $D_{H_2}$ & 20 & 0.196 & 0.1068 & 0.0744 \\
MalNet weighted & $D_F$ / $D_C$ & 20 & 0.714 & 0.0011 & 0.0063 \\
MalNet weighted & $D_F$ / $D_{H_2}$ & 20 & 0.116 & 0.0785 & 0.0638 \\
MalNet weighted & $D_C$ / $D_{H_2}$ & 20 & 0.157 & 0.0834 & 0.0864 \\
COLLAB & $D_F$ / $D_C$ & 5 & 0.714 & 0.0258 & 0.0234 \\
COLLAB & $D_F$ / $D_{H_2}$ & 5 & 0.398 & 0.1058 & 0.0922 \\
COLLAB & $D_C$ / $D_{H_2}$ & 6 & 0.426 & 0.0571 & 0.0288 \\
PROTEINS & $D_F$ / $D_C$ & 16 & 0.764 & 0.0046 & 0.0084 \\
PROTEINS & $D_F$ / $D_{H_2}$ & 14 & 0.570 & 0.1163 & 0.0506 \\
PROTEINS & $D_C$ / $D_{H_2}$ & 13 & 0.568 & 0.1169 & 0.0690 \\
\bottomrule
\end{tabular}
\caption{Hard-partition agreement and cross-distortion transfer at target collision occupancy $H_2(\bar q_G)=2.0$ nats. $n$ counts matched graphs. ARI is invariant to code-label permutation.}
\label{tab:transductive-rate-matched-agreement}
\end{table}

Each cell in the full transfer table averages the accepted graph--target matches across all occupancy targets for that training/evaluation pair.
Each column holds one evaluation fidelity fixed, each row supplies the partition trained using a particular fidelity, and the diagonal is zero by construction.
Coverage can therefore differ between cells.
Comparisons are made within columns because each column is expressed on its evaluation fidelity's scale.

\begin{table}[htbp]
\centering\small
\setlength{\tabcolsep}{4pt}
\begin{tabular}{llrrrr}
\toprule
Source & Training fidelity & Eval. $D_E$ & Eval. $D_F$ & Eval. $D_C$ & Eval. $D_{H_2}$ \\
\midrule
MalNet unweighted & $D_E$ & 0.0000 & -0.0017 & -0.0020 & 0.0732 \\
MalNet unweighted & $D_F$ & 0.0163 & 0.0000 & -0.0024 & 0.0911 \\
MalNet unweighted & $D_C$ & 0.0260 & 0.0058 & 0.0000 & 0.1057 \\
MalNet unweighted & $D_{H_2}$ & 0.0999 & 0.0690 & 0.0738 & 0.0000 \\
\midrule
MalNet weighted & $D_E$ & 0.0000 & 0.0033 & -0.0023 & 0.0658 \\
MalNet weighted & $D_F$ & 0.0131 & 0.0000 & 0.0011 & 0.0785 \\
MalNet weighted & $D_C$ & 0.0149 & 0.0063 & 0.0000 & 0.0834 \\
MalNet weighted & $D_{H_2}$ & 0.1010 & 0.0638 & 0.0864 & 0.0000 \\
\midrule
COLLAB & $D_E$ & 0.0000 & 0.0071 & 0.0509 & 0.2296 \\
COLLAB & $D_F$ & 0.0044 & 0.0000 & 0.0293 & 0.1946 \\
COLLAB & $D_C$ & 0.0381 & 0.0241 & 0.0000 & 0.0597 \\
COLLAB & $D_{H_2}$ & 0.1016 & 0.0746 & 0.0101 & 0.0000 \\
\midrule
PROTEINS & $D_E$ & 0.0000 & -0.0009 & 0.0042 & 0.0758 \\
PROTEINS & $D_F$ & 0.0036 & 0.0000 & 0.0048 & 0.0848 \\
PROTEINS & $D_C$ & 0.0102 & 0.0027 & 0.0000 & 0.0807 \\
PROTEINS & $D_{H_2}$ & 0.0594 & 0.0489 & 0.0577 & 0.0000 \\
\bottomrule
\end{tabular}
\caption{Cross-distortion transfer, averaged over pairwise hard-$H_2$-matched graph--occupancy comparisons. Each entry is excess hard distortion relative to the solution trained for that column's evaluation fidelity. Compare rows within a column.}
\label{tab:transductive-cross-distortion-transfer}
\end{table}

\begin{figure}[htbp]
\centering
\includegraphics[width=0.82\linewidth]{transductive_relational_distortion_partition_examples.png}
\caption{Example partitions at marginal-organization weight 0.50. Rows show a selected PROTEINS graph and a selected COLLAB graph; columns change only the source fidelity optimized by the transductive solver. Positions are fixed within a row, while partition labels and palettes are independently permuted for display. Annotations give each criterion's own hard distortion and hard effective count.}
\label{fig:transductive-relational-distortion-partitions}
\end{figure}

\begin{figure}[htbp]
\centering
\includegraphics[width=0.76\linewidth]{relative_partition_ari_vs_h2.png}
\caption{All six pairwise ARI comparisons in Study 2. The main text highlights three pairs; this figure retains the complete set. Points require at least five graphs meeting the pairwise hard-entropy matching rule.}
\label{fig:transductive-full-rate-matched-ari}
\end{figure}

\begin{figure}[htbp]
\centering
\includegraphics[width=0.76\linewidth]{relative_partition_collision_jaccard_vs_h2.png}
\caption{All six pairwise collision-Jaccard comparisons. Jaccard compares the induced same-code equivalence relations directly, complementing chance-corrected ARI. Points require at least five matched graphs.}
\label{fig:transductive-full-rate-matched-jaccard}
\end{figure}

\section{Inductive Normalized-Cut Study Details}
\label{app:malnet-ncut-details}

MalNet-Tiny retains its official train/validation/test split before cleaning.
Graphs are converted to simple undirected graphs, reduced to their largest connected component, and retained only if that component contains at least 512 nodes.
\Cref{tab:malnet-preprocessing-summary} gives the resulting exclusions.

\begin{table}[htbp]
\centering\small
\begin{tabular}{lrrr}
\toprule
Split & Raw graphs & Retained graphs & Excluded below 512 LCC nodes \\
\midrule
Train & 3500 & 1995 & 1505 \\
Validation & 500 & 275 & 225 \\
Test & 1000 & 584 & 416 \\
\bottomrule
\end{tabular}
\caption{MalNet-Tiny source selection after cleaning.}
\label{tab:malnet-preprocessing-summary}
\end{table}

\begin{table}[htbp]
\centering\small
\begin{tabularx}{\linewidth}{>{\raggedright\arraybackslash}p{0.28\linewidth}X}
\toprule
Component & Setting \\
\midrule
Dataset & MalNet-Tiny, official source splits, labels unused \\
Cleaning & Undirected edge union, self-loop removal, coalesced duplicates, largest connected component \\
Minimum graph size & 512 vertices after LCC extraction \\
Node features & Constant and cleaned undirected degree features, original directed degree features and PageRank, deterministic random-walk return estimates \\
Random-walk estimates & 16 steps, 8 Rademacher probes, split/index-dependent deterministic seeds \\
Assignments & $K=8$ categorical states; $\tau=1.0$ \\
Encoder & Four GraphGPS-style layers, local GIN and global efficient attention, width 128, four heads \\
Additional encoder settings & GELU, no dropout, FFN multiplier 2, graph-context concatenation, linear assignment head \\
Organization weights $\lambda_{\mathrm{org}}$ & $0$, $0.03$, $0.10$, $0.20$, $0.50$ \\
Objective & Minibatch objective $\mathcal L_B$ from \cref{eq:malnet-ncut-loss} \\
Optimizer & AdamW, learning rate $5\times10^{-4}$, zero weight decay \\
Schedule & 50-step warmup from factor 0.1, cosine decay to $10^{-5}$ \\
Training & 250 epochs, batch size 80 \\
Seeds & 1337, 2024, 31415, 27182, 16180 \\
Checkpoint selection & Validation mean hard $\operatorname{Ncut}(G)$ under the fixed-$K$ convention of \cref{eq:hard-ncut} \\
Spectral reference & Sparse normalized-Laplacian embedding, eight smallest eigenvectors, row normalization, deterministic $K$-means discretization \\
Spectral eigensolver & SciPy \texttt{eigsh}, tolerance $10^{-5}$ \\
Spectral $K$-means & Seed 1337, $n_{\mathrm{init}}=8$, maximum 100 iterations, tolerance $10^{-5}$ \\
\bottomrule
\end{tabularx}
\caption{Settings for Study 3.}
\label{tab:malnet-ncut-settings}
\end{table}

The spectral embedding is computed with SciPy's sparse \texttt{eigsh} solver~\citep{virtanen2020scipy}, and its seeded $K$-means discretization uses scikit-learn's \texttt{KMeans} implementation~\citep{pedregosa2011scikit}.
Post-hoc $D_F$ is computed from the selected hard partitions using effective resistances of the cleaned source graph, as in \cref{eq:hard-partition-dirichlet-distortion}.
The spectral reference is computed separately for each retained validation and test graph.

\section{Reconstruction-Trained Flowers102 Study Details}
\label{app:flowers102-details}

All five image reconstruction settings use seed 1337, batch size 64, and 150 training epochs.
AdamW uses learning rate $5\times10^{-4}$, zero weight decay, 100-step warmup from factor 0.1, and cosine decay to $10^{-6}$.
Training applies random resized crops, horizontal flips, and mild color jitter; validation and test use resize followed by deterministic center crop.
The checkpoint criterion is hard-code validation MSE.

The nonstandard split uses the official Flowers102 test images for training, official validation images for validation, and official training images for final test evaluation.
The resulting sizes are 6149, 1020, and 1020, respectively.

\begin{table}[htbp]
\centering\small
\begin{tabularx}{\linewidth}{>{\raggedright\arraybackslash}p{0.34\linewidth}X}
\toprule
Component & Setting \\
\midrule
Input & $256\times256$ RGB, unit-range pixel values \\
Encoder width & 128 hidden channels \\
Residual stack & Two blocks, 64 residual hidden channels \\
Downsampling / upsampling & Three stride-2 analysis stages with GDN and three transposed-convolution synthesis stages, with IGDN after the first two synthesis stages, giving a $32\times32$ latent grid \\
Spatial collision unit & One $b$-bit token per latent position \\
Analysis output & $1\times1$ convolution producing $b$ logits per spatial position \\
Decoder output & Sigmoid RGB reconstruction; no encoder--decoder skips \\
Reconstruction loss & Mean squared error on $[0,1]$ pixels \\
Training objective & $\mathcal L$ from \cref{eq:flowers-objective} \\
Mean-group relaxation & Deterministic signs, separate positive and negative groups per channel, grouped over batch and spatial positions \\
Concentration weight & $\lambda_{\mathrm{conc}}=10$ \\
Threshold-scale EMA decay & 0.99 \\
Threshold-scale floor & $10^{-3}$ \\
Margin gain & $\gamma=3$ \\
Collision pair samples & 8192 minibatch spatial-token pairs \\
Augmentation & Random resized crop, horizontal flip, brightness/hue/saturation/contrast jitter of 0.05 \\
\bottomrule
\end{tabularx}
\caption{Shared settings for Study 4.}
\label{tab:flowers102-shared-settings}
\end{table}
Each collision pair is formed by two independent uniform draws, with replacement, from the minibatch's spatial tokens; self-pairs and same-image pairs are permitted.
These pairs form the minibatch estimate in \cref{eq:flowers-separation} using the threshold collision of \cref{eq:flowers-threshold-collision}.

The mean-group relaxation forms separate positive- and negative-sign groups within each channel over batch and spatial positions, subtracts each nonempty group's mean logit, and adds the resulting centered residual to the corresponding hard sign.
For minibatch size $B$, let $\tilde h_r(x_s,u)$ denote the relaxed latent corresponding to hard sign $h_r(x_s,u)$.
The concentration term is
\begin{equation}
\mathcal L_{\mathrm{conc}}
=\frac{1}{B b H_z W_z}
\sum_{s=1}^{B}\sum_{r=1}^{b}\sum_u
\bigl(\tilde h_r(x_s,u)-h_r(x_s,u)\bigr)^2.
\label{eq:concentration-reduction}
\end{equation}
Here $u$ ranges over the $H_zW_z$ latent spatial positions.
Reported reconstruction metrics use hard signs fed directly to the decoder.

\begin{table}[htbp]
\centering\small
\begin{tabular}{lrrrr}
\toprule
Setting & Bits & $\lambda_{\mathrm{org}}$ & Nominal bpp & Selected epoch \\
\midrule
No organization & 16 & $0$ & 0.25 & 117 \\
Weaker organization & 16 & $10^{-7}$ & 0.25 & 145 \\
Baseline & 16 & $10^{-6}$ & 0.25 & 148 \\
Stronger organization & 16 & $10^{-5}$ & 0.25 & 148 \\
Capacity ablation & 32 & $10^{-6}$ & 0.50 & 150 \\
\bottomrule
\end{tabular}
\caption{Ablations and validation-selected checkpoint epochs.}
\label{tab:flowers102-ablation-settings}
\end{table}

\begin{table}[htbp]
\centering\small
\begin{tabular}{lrrr}
\toprule
Setting & Test MSE & Test PSNR (dB) & Test MS-SSIM \\
\midrule
No organization & 0.01034 & 19.86 & 0.738 \\
Weaker organization & 0.00367 & 24.36 & 0.882 \\
Baseline & 0.00329 & 24.83 & 0.896 \\
Stronger organization & 0.00350 & 24.56 & 0.892 \\
Capacity ablation & 0.00258 & 25.88 & 0.920 \\
\bottomrule
\end{tabular}
\caption{Additional hard-code reconstruction metrics for the same checkpoints as \cref{tab:flowers102-results}. Each setting is a single-seed run.}
\label{tab:flowers102-reconstruction-metrics}
\end{table}

\begin{figure}[htbp]
\centering
\includegraphics[width=\linewidth]{flowers102_reconstruction_examples.png}
\caption{Qualitative Flowers102 test reconstructions. Each triptych shows the original and hard-code reconstructions from the validation-selected 16-bit and 32-bit checkpoints. These examples use the hard sign representation evaluated quantitatively.}
\label{fig:flowers102-reconstruction-examples}
\end{figure}

\section{Teacher-Defined Flowers102 Study Details}
\label{app:flowers102-teacher-details}

The teacher-defined study uses the same source split, image size, and evaluation preprocessing as Study 4, with a single seed of 1337.
Representation learning uses only the baseline-relative teacher relation through favored and disfavored pair requirements.
The secondary synthesis stage begins after the representation checkpoint is selected and keeps the encoder frozen.

\begin{table}[htbp]
\centering\small
\begin{tabularx}{\linewidth}{>{\raggedright\arraybackslash}p{0.29\linewidth}X}
\toprule
Component & Setting \\
\midrule
Dataset split & 6149 train / 1020 validation / 1020 test, as in Study 4 \\
Input & $256\times256$ RGB \\
Teacher & Frozen DINO ViT-S/8, final spatial patch tokens, normalized before cosine evaluation \\
Teacher source & Official \texttt{facebookresearch/dino} \\
Teacher checkpoint & \texttt{dino\_deitsmall8\_pretrain.pth}, official release \\
Teacher grid & $32\times32$ patch-token embeddings \\
Student & The convolutional analysis transform used in Study 4 \\
Hard code & Sixteen zero-threshold sign bits per spatial token \\
Training token sample & 1024 spatial tokens sampled uniformly without replacement per minibatch \\
Evaluation token sample & 1024 spatial tokens sampled uniformly without replacement with fixed validation/test seeds \\
Candidate mask & Cross-image pairs among sampled tokens \\
Temperatures & Code $\tau_c=0.25$; teacher $\tau_T=0.1$ \\
Alignment weight & 1 \\
Representation objective & $\mathcal L_{\mathrm{teacher}}$ from \cref{eq:dino-teacher-loss} \\
Representation optimizer & AdamW, learning rate $5\times10^{-4}$, zero weight decay \\
Representation schedule & 100-step warmup from factor 0.1, cosine decay to $10^{-6}$ \\
Representation training & 150 epochs, batch size 32, seed 1337 \\
Representation selection & Validation teacher loss; selected epoch 146 \\
Decoder stage & Frozen encoder, hard signs only, MSE loss \\
Decoder architecture & The synthesis-transform form used in Study 4 \\
Decoder optimizer & AdamW, learning rate $5\times10^{-4}$, zero weight decay \\
Decoder schedule & 100-step warmup from factor 0.1, cosine decay to $10^{-6}$ \\
Decoder training/selection & 150 epochs, batch size 32; validation MSE selects epoch 140 \\
\bottomrule
\end{tabularx}
\caption{Settings for Study 5.}
\label{tab:flowers102-teacher-settings}
\end{table}

The DINO implementation uses commit \texttt{7c446df5b9f45747937fb0d72314eb9f7b66930a} of the official \texttt{facebookresearch/dino} repository.
The teacher receives the same augmented crop as the student after applying teacher-specific normalization.
The teacher relation and its favored and disfavored weights are defined in \cref{eq:teacher-row-relation,eq:teacher-signed-weights}.

For numerical stability, per-bit agreement probabilities are floored at the floating-point dtype's smallest positive normal value before their logarithms are summed, and the resulting whole-code probability is clipped to $[\varepsilon,1-\varepsilon]$, where $\varepsilon$ is machine epsilon, before evaluating the logit in \cref{eq:teacher-centered-margin}.
For the centered margin to retain a numerically representable region below the neutral collision level, the clipping floor must satisfy $\varepsilon<2^{-b}$; the reported 16-bit loss is evaluated in float32, for which this condition holds.
The clipped loss can differ from the ideal mathematical expression in saturated regions and can have zero gradients through the clipping operation there.
These operations affect only the smooth training loss; reported hard-collision metrics use hard code equality.

The reported favored-pair enrichment and disfavored-pair suppression use pooled empirical token collision as their baseline.
Reported teacher-collision, teacher-cosine, and reconstruction metrics use the conventions in \ref{app:metric-definitions}.

\begin{figure}[htbp]
\centering
\includegraphics[height=0.84\textheight,keepaspectratio]{flowers102_teacher_decoder_reconstructions.png}
\caption{Qualitative Flowers102 test examples from the secondary decoder trained on frozen teacher-guided hard codes. Each pair shows an original image and its reconstruction. Pixel reconstruction did not participate in learning the encoder.}
\label{fig:flowers102-teacher-reconstruction-examples}
\end{figure}

